\documentclass[letterpaper,single,8pt]{article}
\usepackage[margin=2.2cm]{geometry}
\usepackage{xspace}
\usepackage[hidelinks, plainpages=false, pdfpagelabels]{hyperref}

\newcommand{\chap}{Chapter\xspace}
\newcommand{\ie}{i.e.,\xspace}
\newcommand{\eg}{e.g.,\xspace}
\newcommand{\restler}{RESTler\xspace}
\newcommand{\pythia}{Pythia\xspace}
\def\gitlab{GitLab\xspace}


\begin{document}
\title{\vspace{-1.2cm}{\bf Towards Secure and Reliable Software Systems}}
\author{ {\rm Vaggelis Atlidakis} }
\date{}
\maketitle

\begingroup\emergencystretch=1em
Modern software systems are complex, and ensuring their security and reliability is challenging. The widespread adoption of AI coding agents makes this harder by accelerating software development and potentially introducing software errors---including security vulnerabilities and reliability issues---at a scale that traditional testing and security practices cannot keep pace with. Addressing this challenge requires advances along three complementary directions, including \textbf{(i) automatically discovering software errors through testing}, \textbf{(ii) preventing their exploitation through software hardening}, and \textbf{(iii) efficiently remediating known vulnerabilities}. My work over the past decade spans all three directions, with an emphasis on developing practical, scalable solutions to improve the security and reliability of production software systems.
\par\endgroup

\textbf{Automated Software Testing.} A major focus of my work has been designing automated testing systems capable of discovering previously unknown software errors and security vulnerabilities in complex software. During my Ph.D., along with collaborators from Microsoft Research, we introduced the concept of stateful REST API fuzzing and developed \textbf{RESTler}~\cite{restler,restler2}: the first stateful REST API fuzzer. Since 2017, RESTler has been used internally by multiple Microsoft teams for security, reliability, functional correctness, and regression testing, and has found hundreds of bugs in multiple production O365, Azure, and Bing cloud services. In 2020, RESTler was open-sourced, and its usage has been steadily growing since, both inside and outside Microsoft.

Furthermore, I have worked on vulnerability discovery in deep learning frameworks, whose complex software stacks expose low-level native code through high-level programming interfaces. In 2022, along with collaborators from Brown University, we introduced the concept of bottom-up fuzzing and developed \textbf{IvySyn}~\cite{ivysyn}: a fuzzing framework that identifies memory-safety vulnerabilities in native C/C++ kernels and automatically synthesizes high-level proof-of-vulnerability programs demonstrating their reachability through public APIs. IvySyn uncovered 61 previously unknown security vulnerabilities in TensorFlow and PyTorch, resulting in 39 CVEs. Building on this direction, our recent work, \textbf{FuzzWarp}~\cite{fuzzwarp}, extends bottom-up fuzzing to GPU kernels in deep learning frameworks, introducing GPU-aware mechanisms for precise device dispatch, input mutation, and memory-error detection. FuzzWarp outperformed state-of-the-art DL fuzzers, uncovering 99 previously unknown GPU-specific memory-safety bugs in TensorFlow and PyTorch.

\textbf{Software Systems Hardening.} Along with collaborators from Brown University, we have developed practical hardening mechanisms that leverage compiler-assisted instrumentation, operating system abstractions, and hardware security features. First, we developed \textbf{SysXCHG}~\cite{sysxchg}, a compiler-assisted mechanism for enforcing the principle of least privilege through adaptive system-call filtering, allowing applications to dynamically restrict their access to the Linux kernel. Second, we introduced \textbf{FineIBT}~\cite{fineibt}, a fine-grained control-flow integrity mechanism that builds on Intel's Indirect Branch Tracking technology to protect applications against control-flow hijacking attacks. FineIBT was adopted by Intel, and a variant was subsequently upstreamed to the Linux kernel. Third, we developed \textbf{Eclipse}~\cite{eclipse}, a compiler-assisted defense against speculative memory-error attacks, preventing attackers from combining memory corruption vulnerabilities with speculative execution to bypass existing security mitigations.

\textbf{Vulnerability Remediation.} Beyond academic research, I have also held software engineering positions in industry, where I have developed and shipped practical security solutions for production systems. At Lacework, I was part of the Code Security team led by Peter O'Hearn, where I worked with Patrice Godefroid on \textbf{Code-Aware Agent (CAA)}~\cite{caa}: a system for vulnerability risk management. Our key insight was that not all vulnerable software packages are actually executed in production, and therefore not all vulnerabilities pose the same immediate security risk. To this end, CAA continuously monitors software package activity at runtime, allowing security teams to prioritize vulnerabilities in actively executed code. Internal anecdotal studies suggested that this approach, when used by our customers, reduced the total number of vulnerable packages requiring immediate attention by almost 80\% daily.

\enlargethispage{2\baselineskip}
Overall, my goal is to continue developing practical, scalable technologies that improve the security and reliability of modern software systems. I believe that my background in automated software testing, systems hardening, and vulnerability risk management, combined with my experience in both academic research and industrial software engineering, aligns closely with Bloomberg's commitment to technological innovation. I would welcome the opportunity to collaborate with Bloomberg's organization on advancing the state of the art in software security and reliability, and translating research ideas into practical solutions.

\pagebreak
\bibliographystyle{abbrv}
\bibliography{statement}


\end{document}
